% from ega4-1-fr.tex
\oldpage[IV-1]{242}

\begin{env}[1.9.2.ii]
\label{IV.1.9.2.ii-fr}
\emph{On a}
\[
  u(X)=\bigcap_{\alpha\in I}u_\alpha(X_\alpha). \tag{1.9.2.1}
  \label{IV.1.9.2.1-fr}
\]
\end{env}

L'assertion (i) n'est autre que la traduction de (1.9.1). Pour démontrer
(ii), notons que, puisque $u$ se factorise en
$X\to X_\alpha\xrightarrow{u_\alpha}Y$, le premier membre de (1.9.2.1) est
contenu dans le second. Inversement, soit $y\in Y-u(X)$~; on a
$u^{-1}(y)=\varnothing$, et $u^{-1}(y)$ est l'espace sous-jacent au
$k(y)$-préschéma $\operatorname{Spec}(A\otimes_Bk(y))$~; comme dans la
catégorie des $B$-modules le foncteur $\varinjlim$ commute au produit
tensoriel, $A\otimes_Bk(y)$ est limite inductive du système inductif d'anneaux
$A_\alpha\otimes_Bk(y)$~; le lemme (1.9.1) montre qu'il existe $\alpha$ tel
que $A_\alpha\otimes_Bk(y)=0$, c'est-à-dire
$u_\alpha^{-1}(y)=\varnothing$, donc $y\notin u_\alpha(X_\alpha)$. C.Q.F.D.

\begin{proposition}[1.9.3]
\label{IV.1.9.3-fr}
Soient $X$ un préschéma quasi-compact et quasi-séparé, $E$ une partie de $X$,
$(X_\alpha)_{\alpha\in I}$ un recouvrement ouvert affine de $X$, et pour
tout $\alpha\in I$, posons $E_\alpha=E\cap X_\alpha$. Les conditions
suivantes sont équivalentes~:
\begin{enumerate}
  \item[a)] Il existe un morphisme quasi-compact $f:X'\to X$ tel que
  $E=f(X')$.
  \item[a$'$)] Pour tout $\alpha$, il existe un morphisme quasi-compact
  $f_\alpha:X'_\alpha\to X_\alpha$ tel que
  $E_\alpha=f_\alpha(X'_\alpha)$.
  \item[a$''$)] Pour tout $\alpha$, il existe un schéma affine $X''_\alpha$
  et un morphisme $f'_\alpha:X''_\alpha\to X_\alpha$ tel que
  $f'_\alpha(X''_\alpha)=E_\alpha$.
  \item[b)] $E$ est intersection de parties constructibles de $X$.
  \item[b$'$)] $F=X-E$ est réunion de parties constructibles de $X$.
\end{enumerate}
\end{proposition}

\begin{proof}
Il est clair que $b)$ et $b')$ sont équivalentes. La condition $a)$ entraîne
$a')$ en prenant pour $X'_\alpha$ le préschéma induit sur l'ouvert
$f^{-1}(X_\alpha)$ de $X'$ et pour $f_\alpha$ la restriction de $f$. Pour
voir que $a')$ entraîne $a'')$, il suffit de remarquer que $X'_\alpha$ est
réunion d'ouverts affines $X''_{\alpha\lambda}$ ($\lambda\in L$)~; soit
$X''_\alpha$ le préschéma \emph{somme} des $X''_{\alpha\lambda}$
($\lambda\in L$), qui est un schéma affine~; si
$g_\alpha:X''_\alpha\to X'_\alpha$ est le morphisme coïncidant avec
l'identité dans chacun des $X''_{\alpha\lambda}$, il est clair que si l'on
pose $f'_\alpha=f_\alpha\circ g_\alpha$, on a
$E_\alpha=f'_\alpha(X''_\alpha)$ puisque $g_\alpha$ est surjectif. Pour
montrer que $a'')$ entraîne $a)$, notons qu'il y a une partie finie $J$ de
$I$ telle que les $X_\alpha$ d'indices $\alpha\in J$ recouvrent $X$~; soit
$X'$ le préschéma somme des $X''_\alpha$ pour $\alpha\in J$, et soit
$f:X'\to X$ le morphisme coïncidant avec $j_\alpha\circ f'_\alpha$ pour tout
$\alpha\in J$, où $j_\alpha:X_\alpha\to X$ est l'injection canonique. Comme
$E$ est réunion des $E\cap X_\alpha$ pour $\alpha\in J$, on a bien
$E=f(X')$ et il reste à voir que $f$ est un morphisme quasi-compact~; mais
l'hypothèse que $X$ est quasi-séparé entraîne que $j_\alpha$ est
quasi-compact (1.2.7), donc il en est de même de
$j_\alpha\circ f'_\alpha$ (1.1.1 et 1.1.2) et finalement de $f$ (1.1.6).

Reste donc à prouver l'équivalence de $a'')$ et de $b)$. Prouvons en premier
lieu que $a'')$ entraîne $b)$~: considérons une partie finie $J$ de $I$ telle
que les $X_\alpha$ pour $\alpha\in J$ recouvrent $X$~; il suffira de montrer
que, pour tout $\alpha\in J$, $E_\alpha$ est intersection de parties
constructibles $F_{\alpha\lambda}$ de $X_\alpha$ ($\lambda\in L_\alpha$)~;
en effet, tout $F_{\alpha\lambda}$ est aussi une partie constructible de $X$
en vertu de (0\textsubscript{III}, 9.1.8, (ii)), car l'hypothèse que $X$ est
quasi-séparé entraîne que tout $X_\alpha$ est rétrocompact dans $X$ (1.2.7)~;
$E$ étant réunion des $E_\alpha$ pour $\alpha\in J$, est intersection des
réunions (finies) $\bigcup_{\alpha\in J}F_{\alpha,\lambda(\alpha)}$, pour
tous les choix des $\lambda(\alpha)\in L_\alpha$~; ces réunions étant
constructibles dans $X$, cela établit notre assertion. On peut donc se borner
au cas

